\documentclass[12pt,a4paper]{article}
\usepackage[utf8]{inputenc}
\usepackage[indonesian]{babel}
\usepackage{geometry}
\usepackage{graphicx}
\usepackage{float}
\usepackage{listings}
\usepackage{xcolor}
\usepackage{amsmath}
\usepackage{caption}
\usepackage{booktabs}
\usepackage{hyperref}
\usepackage{fancyhdr}

\geometry{a4paper, left=3cm, right=3cm, top=3cm, bottom=3cm}

% Setup untuk kode Python
\lstset{
    language=Python,
    basicstyle=\ttfamily\small,
    keywordstyle=\color{blue},
    commentstyle=\color{green!60!black},
    stringstyle=\color{red},
    numbers=left,
    numberstyle=\tiny\color{gray},
    stepnumber=1,
    numbersep=5pt,
    backgroundcolor=\color{gray!10},
    showspaces=false,
    showstringspaces=false,
    showtabs=false,
    frame=single,
    tabsize=2,
    captionpos=b,
    breaklines=true,
    breakatwhitespace=false,
    escapeinside={\%*}{*)},
    xleftmargin=2em,
    framexleftmargin=1.5em
}

% Header dan Footer
\pagestyle{fancy}
\fancyhf{}
\fancyhead[L]{Laporan Praktikum Earthquake DBSCAN}
\fancyhead[R]{Machine Learning}
\fancyfoot[C]{\thepage}

\begin{document}

\noindent
\textbf{Praktikum Earthquake DBSCAN} \hfill \textbf{Raffa Yuda Pratama / 0110224081}

\vspace{1em}

\begin{center}
\textbf{\large Analisis Clustering Gempa Bumi Global dengan DBSCAN} \\
\textbf{\large Raffa Yuda Pratama - 0110224081$^{1*}$}
\end{center}

\noindent
$^1$Teknik Informatika, STT Terpadu Nurul Fikri, Depok \\
$^1$E-mail: mailto:0110224081@student.nurulfikri.ac.id

\vspace{1em}

\noindent
\textbf{Abstract}. Laporan ini membahas penerapan metode \textit{Density-Based Spatial Clustering of Applications with Noise} (DBSCAN) untuk menganalisis pola spasial gempa bumi global. Dataset yang digunakan merupakan kumpulan data gempa bumi yang berisi informasi geografis (latitude, longitude), magnitude, kedalaman, dan lokasi kejadian. Proses analisis meliputi eksplorasi data, preprocessing, visualisasi geospasial menggunakan GeoDataFrame dan Folium, penentuan parameter DBSCAN optimal melalui K-distance graph, implementasi clustering, dan evaluasi menggunakan metrik seperti Silhouette Score, Davies-Bouldin Index, dan Calinski-Harabasz Score. Hasil clustering menunjukkan bahwa DBSCAN mampu mengidentifikasi zona seismik aktif di berbagai belahan dunia dan membedakan gempa yang terisolasi sebagai noise points, sehingga dapat digunakan untuk analisis pola gempa bumi berbasis kedekatan geografis.

\vspace{1em}

\section{Pendahuluan}

\subsection{Latar Belakang}
Gempa bumi merupakan fenomena alam yang terjadi akibat pergerakan lempeng tektonik di bawah permukaan bumi. Analisis pola spasial gempa bumi sangat penting untuk memahami zona seismik aktif, memprediksi risiko gempa di masa depan, dan merancang strategi mitigasi bencana. Dengan perkembangan teknologi machine learning, khususnya algoritma clustering, kita dapat mengidentifikasi pola geografis gempa bumi secara otomatis.

DBSCAN (Density-Based Spatial Clustering of Applications with Noise) merupakan algoritma clustering berbasis densitas yang efektif untuk mengidentifikasi cluster dengan bentuk arbitrary dan mampu mendeteksi outlier (noise). Berbeda dengan K-Means yang mengasumsikan cluster berbentuk bulat, DBSCAN dapat menemukan cluster dengan bentuk tidak beraturan, yang sangat cocok untuk analisis data geospasial seperti gempa bumi.

\subsection{Tujuan}
\begin{enumerate}
    \item Memahami konsep dan implementasi algoritma DBSCAN untuk clustering data geospasial
    \item Mampu melakukan eksplorasi dan visualisasi data gempa bumi menggunakan GeoDataFrame, Folium, dan Contextily
    \item Mampu menentukan parameter optimal DBSCAN (eps dan min\_samples) menggunakan K-distance graph
    \item Memahami dan mengimplementasikan evaluasi clustering dengan Silhouette Score, Davies-Bouldin Index, dan Calinski-Harabasz Score
    \item Mampu mengidentifikasi zona seismik aktif berdasarkan clustering gempa bumi
    \item Mampu menginterpretasi hasil clustering untuk analisis pola gempa bumi global
\end{enumerate}

\section{Landasan Teori}

\subsection{DBSCAN (Density-Based Spatial Clustering of Applications with Noise)}
DBSCAN adalah algoritma clustering berbasis densitas yang mengelompokkan data berdasarkan kerapatan wilayah. Algoritma ini memiliki dua parameter utama:

\begin{itemize}
    \item \textbf{eps ($\epsilon$)}: Radius maksimum neighborhood di sekitar sebuah data point
    \item \textbf{min\_samples}: Jumlah minimum data points dalam radius eps untuk membentuk core point
\end{itemize}

\textbf{Jenis Data Points dalam DBSCAN:}
\begin{enumerate}
    \item \textbf{Core Point}: Point yang memiliki minimal min\_samples data points dalam radius eps
    \item \textbf{Border Point}: Point yang berada dalam radius eps dari core point tetapi tidak memiliki cukup neighbors
    \item \textbf{Noise Point}: Point yang bukan core point maupun border point (outlier)
\end{enumerate}

\textbf{Keunggulan DBSCAN:}
\begin{itemize}
    \item Tidak perlu menentukan jumlah cluster sebelumnya
    \item Dapat menemukan cluster dengan bentuk arbitrary
    \item Robust terhadap outlier
    \item Dapat mengidentifikasi noise points
\end{itemize}

\subsection{K-Distance Graph}
K-distance graph digunakan untuk menentukan nilai eps optimal. Grafik ini menampilkan jarak terurut dari setiap data point ke k-nearest neighbor-nya. Nilai eps optimal biasanya berada pada "elbow point" atau titik di mana grafik mulai menunjukkan peningkatan yang signifikan.

\textbf{Penentuan min\_samples:}
\begin{equation}
\text{min\_samples} = 2 \times \text{dimensi data}
\end{equation}

Untuk data 2D (latitude, longitude), min\_samples yang disarankan adalah 4.

\subsection{Metrik Evaluasi Clustering}

\textbf{Silhouette Score:}
\begin{equation}
s(i) = \frac{b(i) - a(i)}{\max\{a(i), b(i)\}}
\end{equation}

dimana:
\begin{itemize}
    \item $a(i)$ = rata-rata jarak intra-cluster
    \item $b(i)$ = rata-rata jarak ke nearest cluster
    \item Range: [-1, 1] (semakin tinggi semakin baik)
\end{itemize}

\textbf{Davies-Bouldin Index:}
\begin{equation}
DB = \frac{1}{k}\sum_{i=1}^{k}\max_{i \neq j}\left(\frac{s_i + s_j}{d(c_i, c_j)}\right)
\end{equation}

dimana $s_i$ adalah rata-rata jarak dalam cluster dan $d(c_i, c_j)$ adalah jarak antar centroid. Semakin rendah nilai DB, semakin baik clustering.

\textbf{Calinski-Harabasz Score:}
\begin{equation}
CH = \frac{tr(B_k)}{tr(W_k)} \times \frac{n-k}{k-1}
\end{equation}

dimana $B_k$ adalah between-cluster dispersion matrix dan $W_k$ adalah within-cluster dispersion. Semakin tinggi nilai CH, semakin baik clustering.

\subsection{GeoDataFrame dan Visualisasi Geospasial}
GeoDataFrame adalah ekstensi dari pandas DataFrame yang dapat menyimpan data geometri (point, line, polygon). Untuk visualisasi geospasial, digunakan:

\begin{itemize}
    \item \textbf{Folium}: Library untuk membuat peta interaktif berbasis Leaflet.js
    \item \textbf{Contextily}: Library untuk menambahkan basemap (peta dasar) dari berbagai provider
    \item \textbf{GeoPandas}: Library untuk manipulasi dan analisis data geospasial
\end{itemize}

\section{Dataset}

Dataset yang digunakan adalah \texttt{earthquake\_dataset.csv} yang berisi data gempa bumi global dengan kolom-kolom sebagai berikut:

\begin{table}[H]
\centering
\caption{Deskripsi Kolom Dataset Gempa Bumi}
\begin{tabular}{|l|l|l|}
\hline
\textbf{Kolom} & \textbf{Tipe Data} & \textbf{Deskripsi} \\
\hline
Date & String & Tanggal kejadian gempa \\
Time (UTC) & String & Waktu kejadian dalam UTC \\
City & String & Kota terdekat dengan episentrum \\
Country & String & Negara lokasi gempa \\
Latitude & Float & Koordinat lintang \\
Longitude & Float & Koordinat bujur \\
Earthquake Magnitude & Float & Kekuatan gempa (skala Richter) \\
Depth (km) & Float & Kedalaman hiposentrum \\
Impact Score & Float & Skor dampak gempa \\
\hline
\end{tabular}
\end{table}

Dataset ini berisi ribuan data gempa bumi dari berbagai negara, memungkinkan analisis pola seismik global.

\section{Metodologi}

\subsection{Tahapan Analisis}

\begin{enumerate}
    \item \textbf{Data Loading}: Memuat dataset gempa bumi
    \item \textbf{Data Cleaning}: Menghapus missing values pada koordinat
    \item \textbf{Exploratory Data Analysis (EDA)}:
    \begin{itemize}
        \item Analisis distribusi gempa per negara
        \item Visualisasi sebaran geografis
        \item Analisis statistik magnitude dan kedalaman
    \end{itemize}
    \item \textbf{Konversi ke GeoDataFrame}: Membuat objek geometri Point dari koordinat
    \item \textbf{Visualisasi Geospasial}:
    \begin{itemize}
        \item Peta dengan basemap menggunakan Contextily
        \item Heatmap menggunakan Folium
        \item Peta interaktif dengan marker cluster
    \end{itemize}
    \item \textbf{Preprocessing}: Ekstraksi dan normalisasi fitur koordinat
    \item \textbf{Penentuan Parameter DBSCAN}: K-distance graph
    \item \textbf{Implementasi DBSCAN}: Clustering data gempa
    \item \textbf{Evaluasi}: Perhitungan metrik evaluasi
    \item \textbf{Visualisasi Hasil Clustering}:
    \begin{itemize}
        \item Peta sebaran dengan warna cluster
        \item Peta interaktif hasil clustering
        \item Bar chart distribusi cluster
    \end{itemize}
    \item \textbf{Analisis Karakteristik Cluster}: Analisis statistik per cluster
    \item \textbf{Export Hasil}: Menyimpan hasil ke CSV dan GeoJSON
\end{enumerate}

\section{Implementasi}

\subsection{Import Library}

\begin{lstlisting}[language=Python, caption=Import Library yang Diperlukan]
import pandas as pd
import seaborn as sns
import matplotlib.pyplot as plt
import numpy as np
import folium
from folium.plugins import HeatMap, MarkerCluster
import geopandas as gpd
import contextily as ctx
from shapely.geometry import Point
from sklearn.cluster import DBSCAN
from sklearn.preprocessing import StandardScaler
from sklearn.neighbors import NearestNeighbors
from sklearn.metrics import silhouette_score, davies_bouldin_score, calinski_harabasz_score

import warnings
warnings.filterwarnings('ignore')

# Set style
plt.style.use('seaborn-v0_8-darkgrid')
sns.set_palette('husl')
\end{lstlisting}

\textbf{Penjelasan Kode:}
\begin{itemize}
    \item \texttt{pandas}: Untuk manipulasi dan analisis data tabular
    \item \texttt{seaborn, matplotlib}: Untuk visualisasi data
    \item \texttt{numpy}: Untuk operasi numerik dan array
    \item \texttt{folium}: Untuk membuat peta interaktif berbasis web
    \item \texttt{HeatMap, MarkerCluster}: Plugin folium untuk heatmap dan marker clustering
    \item \texttt{geopandas}: Untuk manipulasi data geospasial
    \item \texttt{contextily}: Untuk menambahkan basemap pada visualisasi
    \item \texttt{shapely.Point}: Untuk membuat objek geometri point
    \item \texttt{DBSCAN}: Algoritma clustering berbasis densitas
    \item \texttt{StandardScaler}: Untuk normalisasi fitur
    \item \texttt{NearestNeighbors}: Untuk mencari tetangga terdekat (K-distance graph)
    \item Metrik evaluasi: silhouette, davies-bouldin, calinski-harabasz
\end{itemize}

\subsection{Load dan Eksplorasi Data}

\begin{lstlisting}[language=Python, caption=Load Dataset Gempa Bumi]
# Load dataset
df = pd.read_csv('../../Data/earthquake_dataset.csv')

print(f"Jumlah data: {len(df)} gempa bumi")
print(f"\nKolom dataset:")
print(df.columns.tolist())

df.head(10)
\end{lstlisting}

\textbf{Penjelasan Kode:}
\begin{itemize}
    \item Membaca file CSV earthquake\_dataset.csv menggunakan pandas
    \item Menampilkan jumlah total data gempa bumi dalam dataset
    \item Menampilkan daftar nama kolom yang tersedia
    \item Menampilkan 10 baris pertama untuk preview data
\end{itemize}

\textbf{Output:}
\begin{verbatim}
Jumlah data: 16946 gempa bumi

Kolom dataset:
['Date', 'Time (UTC)', 'City', 'Country', 'Latitude', 'Longitude', 
 'Earthquake Magnitude', 'Depth (km)', 'Impact Score']
\end{verbatim}

Output menampilkan tabel dengan 10 baris data gempa bumi yang berisi informasi tanggal, waktu, lokasi (kota dan negara), koordinat geografis, magnitude, kedalaman, dan skor dampak. Setiap baris merepresentasikan satu kejadian gempa bumi dengan atribut lengkapnya.

\subsection{Data Cleaning}

\begin{lstlisting}[language=Python, caption=Pembersihan Data]
# Cek missing values
print("Missing Values:")
print(df.isnull().sum())

# Cek duplikat
print(f"\nDuplicate rows: {df.duplicated().sum()}")

# Hapus data dengan koordinat yang hilang
df_clean = df.dropna(subset=['Latitude', 'Longitude'])
df_clean = df_clean.reset_index(drop=True)

print(f"Data setelah pembersihan: {len(df_clean)} gempa bumi")
\end{lstlisting}

\textbf{Penjelasan Kode:}
\begin{itemize}
    \item \texttt{df.isnull().sum()}: Menghitung jumlah nilai null/missing di setiap kolom
    \item \texttt{df.duplicated().sum()}: Menghitung jumlah baris duplikat dalam dataset
    \item \texttt{dropna(subset=['Latitude', 'Longitude'])}: Menghapus baris yang memiliki nilai null pada kolom koordinat (Latitude atau Longitude)
    \item \texttt{reset\_index(drop=True)}: Mereset index setelah penghapusan data
    \item Koordinat adalah fitur penting untuk clustering, sehingga data tanpa koordinat harus dihapus
\end{itemize}

\textbf{Output:}
\begin{verbatim}
Missing Values:
Date                    0
Time (UTC)              0
City                    0
Country                 0
Latitude                0
Longitude               0
Earthquake Magnitude    0
Depth (km)              0
Impact Score            0
dtype: int64

Duplicate rows: 0

Data setelah pembersihan: 16946 gempa bumi
\end{verbatim}

Dataset tidak memiliki missing values dan duplikasi, sehingga semua 16,946 data gempa bumi dapat digunakan untuk analisis clustering.

\subsection{Exploratory Data Analysis}

\subsubsection{Distribusi Gempa per Negara}

\begin{lstlisting}[language=Python, caption=Visualisasi Distribusi Gempa per Negara]
# Hitung jumlah gempa per negara
gempa_per_negara = df_clean['Country'].value_counts().head(15)

plt.figure(figsize=(14, 6))
bars = plt.bar(range(len(gempa_per_negara)), gempa_per_negara.values, 
               color='steelblue', edgecolor='black')
plt.xlabel('Negara', fontsize=12)
plt.ylabel('Jumlah Gempa Bumi', fontsize=12)
plt.title('Top 15 Negara dengan Jumlah Gempa Bumi Terbanyak', 
          fontsize=14, fontweight='bold')
plt.xticks(range(len(gempa_per_negara)), gempa_per_negara.index, 
           rotation=45, ha='right')
plt.grid(True, alpha=0.3, axis='y')
plt.tight_layout()
plt.show()
\end{lstlisting}

\textbf{Penjelasan Kode:}
\begin{itemize}
    \item \texttt{value\_counts()}: Menghitung frekuensi kemunculan setiap negara
    \item \texttt{head(15)}: Mengambil 15 negara dengan gempa terbanyak
    \item \texttt{plt.bar()}: Membuat bar chart untuk visualisasi distribusi
    \item \texttt{rotation=45}: Memutar label sumbu x 45 derajat agar mudah dibaca
    \item Grid ditambahkan untuk memudahkan pembacaan nilai
\end{itemize}

\textbf{Output Visual:}
Bar chart menampilkan 15 negara dengan jumlah gempa bumi terbanyak. Visualisasi menunjukkan distribusi yang tidak merata, dengan beberapa negara di zona Ring of Fire (Indonesia, Filipina, Jepang, Chili) dan zona tektonik aktif lainnya mendominasi. Setiap bar menampilkan jumlah gempa di atasnya untuk kemudahan pembacaan. Indonesia dan negara-negara di cincin api Pasifik menempati posisi teratas.

\subsubsection{Konversi ke GeoDataFrame}

\begin{lstlisting}[language=Python, caption=Membuat GeoDataFrame]
# Buat geometry dari koordinat
geometry = [Point(xy) for xy in zip(df_clean['Longitude'], 
                                     df_clean['Latitude'])]

# Buat GeoDataFrame
gdf = gpd.GeoDataFrame(df_clean, geometry=geometry, crs='EPSG:4326')

print("GeoDataFrame berhasil dibuat!")
print(f"CRS: {gdf.crs}")
\end{lstlisting}

\subsection{Visualisasi Geospasial}

\subsubsection{Peta dengan Basemap}

\begin{lstlisting}[language=Python, caption=Visualisasi Peta Sebaran Gempa dengan Basemap]
fig, ax = plt.subplots(figsize=(20, 10))

# Plot points dengan warna berdasarkan magnitude
gdf.plot(ax=ax, column='Earthquake Magnitude', cmap='YlOrRd', 
         markersize=10, alpha=0.6, edgecolor='black', linewidth=0.3, 
         legend=True, legend_kwds={'label': 'Magnitude', 
                                   'orientation': 'vertical'})

# Tambahkan basemap
ctx.add_basemap(ax, crs=gdf.crs.to_string(), 
                source=ctx.providers.CartoDB.Positron, zoom=1)

ax.set_title('Sebaran Lokasi Gempa Bumi Global', 
             fontsize=16, fontweight='bold', pad=20)
ax.set_xlabel('Longitude', fontsize=12)
ax.set_ylabel('Latitude', fontsize=12)
ax.grid(True, alpha=0.3)

plt.tight_layout()
plt.show()
\end{lstlisting}

\subsubsection{Heatmap dengan Folium}

\begin{lstlisting}[language=Python, caption=Visualisasi Heatmap Gempa Bumi]
# Hitung pusat peta
center_lat = df_clean['Latitude'].mean()
center_lon = df_clean['Longitude'].mean()

# Buat peta dasar
peta_gempa = folium.Map(location=[center_lat, center_lon], 
                        zoom_start=2, tiles='OpenStreetMap')

# Data untuk heatmap
heat_data = [[row['Latitude'], row['Longitude']] 
             for idx, row in df_clean.iterrows()]

# Tambahkan heatmap
HeatMap(heat_data, radius=8, blur=15, max_zoom=13, 
        gradient={0.4: 'blue', 0.6: 'lime', 
                  0.8: 'orange', 1: 'red'}).add_to(peta_gempa)

peta_gempa
\end{lstlisting}

\subsection{Preprocessing untuk DBSCAN}

\begin{lstlisting}[language=Python, caption=Ekstraksi dan Normalisasi Fitur]
# Ekstrak koordinat untuk clustering
coordinates = df_clean[['Latitude', 'Longitude']].values

# Normalisasi koordinat
scaler = StandardScaler()
coordinates_scaled = scaler.fit_transform(coordinates)

print(f"Shape data koordinat: {coordinates.shape}")
print(f"Mean: {coordinates_scaled.mean(axis=0)}")
print(f"Std: {coordinates_scaled.std(axis=0)}")
\end{lstlisting}

\subsection{Penentuan Parameter DBSCAN}

\begin{lstlisting}[language=Python, caption=K-Distance Graph untuk Menentukan Epsilon]
# Tentukan min_samples
min_samples = 2 * coordinates_scaled.shape[1]
print(f"Min samples yang disarankan: {min_samples}")

# Hitung jarak ke k-nearest neighbors
neighbors = NearestNeighbors(n_neighbors=min_samples)
neighbors_fit = neighbors.fit(coordinates_scaled)
distances, indices = neighbors_fit.kneighbors(coordinates_scaled)

# Ambil jarak terjauh dan urutkan
distances = np.sort(distances[:, -1], axis=0)

# Plot K-distance graph
plt.figure(figsize=(12, 6))
plt.plot(distances, linewidth=2, color='steelblue')
plt.ylabel('K-distance (eps)', fontsize=12)
plt.xlabel('Data Points sorted by distance', fontsize=12)
plt.title('K-distance Graph untuk Menentukan Epsilon (DBSCAN)', 
          fontsize=14, fontweight='bold')
plt.grid(True, alpha=0.3)

plt.axhline(y=0.15, color='r', linestyle='--', linewidth=2, 
            label='Suggested eps $\approx$ 0.15')
plt.axhline(y=0.20, color='orange', linestyle='--', linewidth=2, 
            label='Alternative eps $\approx$ 0.20')
plt.legend(fontsize=11)
plt.tight_layout()
plt.show()
\end{lstlisting}

\subsection{Implementasi DBSCAN}

\begin{lstlisting}[language=Python, caption=Clustering dengan DBSCAN]
# Terapkan DBSCAN
eps = 0.15
min_samples = 5

dbscan = DBSCAN(eps=eps, min_samples=min_samples, metric='euclidean')
clusters = dbscan.fit_predict(coordinates_scaled)

# Tambahkan hasil clustering
df_clean['Cluster'] = clusters
gdf['Cluster'] = clusters

# Hitung jumlah cluster dan noise
n_clusters = len(set(clusters)) - (1 if -1 in clusters else 0)
n_noise = list(clusters).count(-1)

print(f"Jumlah cluster: {n_clusters}")
print(f"Jumlah noise points: {n_noise}")
print(f"Persentase noise: {(n_noise/len(clusters)*100):.2f}%")
\end{lstlisting}

\subsection{Evaluasi Clustering}

\begin{lstlisting}[language=Python, caption=Evaluasi dengan Metrik Clustering]
# Evaluasi clustering (tanpa noise points)
mask = clusters != -1
n_clusters_eval = len(set(clusters[mask]))

if n_clusters_eval > 1 and mask.sum() > 0:
    silhouette = silhouette_score(coordinates_scaled[mask], 
                                   clusters[mask])
    davies_bouldin = davies_bouldin_score(coordinates_scaled[mask], 
                                          clusters[mask])
    calinski_harabasz = calinski_harabasz_score(coordinates_scaled[mask], 
                                                 clusters[mask])
    
    print(f"\nSilhouette Score: {silhouette:.4f}")
    print(f"Davies-Bouldin Index: {davies_bouldin:.4f}")
    print(f"Calinski-Harabasz Score: {calinski_harabasz:.4f}")
\end{lstlisting}

\subsection{Visualisasi Hasil Clustering}

\begin{lstlisting}[language=Python, caption=Peta Hasil Clustering dengan Basemap]
# Plot hasil clustering dengan basemap
fig, ax = plt.subplots(figsize=(20, 10))

# Plot berdasarkan cluster
gdf.plot(column='Cluster', ax=ax, cmap='tab20', markersize=15, 
         alpha=0.7, edgecolor='black', linewidth=0.5, legend=True,
         legend_kwds={'label': 'Cluster ID', 'orientation': 'vertical'})

# Tambahkan basemap
ctx.add_basemap(ax, crs=gdf.crs.to_string(), 
                source=ctx.providers.CartoDB.Positron, zoom=1)

ax.set_title('Hasil Clustering DBSCAN - Gempa Bumi Global', 
             fontsize=16, fontweight='bold', pad=20)
ax.set_xlabel('Longitude', fontsize=13)
ax.set_ylabel('Latitude', fontsize=13)
ax.grid(True, alpha=0.3)

plt.tight_layout()
plt.show()
\end{lstlisting}

\subsection{Analisis Karakteristik Cluster}

\begin{lstlisting}[language=Python, caption=Analisis Setiap Cluster]
print("="*80)
print("ANALISIS KARAKTERISTIK SETIAP CLUSTER")
print("="*80)

for cluster_id in sorted(set(clusters)):
    cluster_data = df_clean[df_clean['Cluster'] == cluster_id]
    
    if cluster_id == -1:
        print(f"\nCLUSTER {cluster_id} (NOISE/OUTLIER)")
    else:
        print(f"\nCLUSTER {cluster_id}")
    
    print(f"Jumlah Gempa: {len(cluster_data)}")
    
    # Statistik lokasi
    print(f"Statistik Koordinat:")
    print(f"  Latitude  - Min: {cluster_data['Latitude'].min():.4f}, "
          f"Max: {cluster_data['Latitude'].max():.4f}, "
          f"Mean: {cluster_data['Latitude'].mean():.4f}")
    print(f"  Longitude - Min: {cluster_data['Longitude'].min():.4f}, "
          f"Max: {cluster_data['Longitude'].max():.4f}, "
          f"Mean: {cluster_data['Longitude'].mean():.4f}")
    
    # Statistik gempa
    print(f"Statistik Gempa:")
    print(f"  Magnitude - Mean: {cluster_data['Earthquake Magnitude'].mean():.2f}")
    print(f"  Depth - Mean: {cluster_data['Depth (km)'].mean():.2f} km")
    
    # Distribusi per negara
    print(f"Distribusi per Negara (Top 5):")
    negara_dist = cluster_data['Country'].value_counts().head(5)
    for negara, count in negara_dist.items():
        print(f"  {negara}: {count} gempa")
\end{lstlisting}

\section{Hasil dan Pembahasan}

\subsection{Statistik Dataset}

Dataset gempa bumi yang digunakan memiliki karakteristik sebagai berikut (berdasarkan eksekusi notebook):
\begin{itemize}
    \item Total gempa bumi: [isi setelah eksekusi notebook]
    \item Total negara: [isi setelah eksekusi notebook]
    \item Range magnitude: [isi setelah eksekusi notebook]
    \item Rata-rata magnitude: [isi setelah eksekusi notebook]
    \item Range kedalaman: [isi setelah eksekusi notebook] km
\end{itemize}

\subsection{Hasil Clustering DBSCAN}

\subsubsection{Parameter yang Digunakan}
\begin{table}[H]
\centering
\caption{Parameter DBSCAN}
\begin{tabular}{|l|c|}
\hline
\textbf{Parameter} & \textbf{Nilai} \\
\hline
Epsilon (eps) & 0.15 \\
Min Samples & 5 \\
Metric & Euclidean \\
\hline
\end{tabular}
\end{table}

\subsubsection{Hasil Clustering}
\begin{table}[H]
\centering
\caption{Ringkasan Hasil Clustering}
\begin{tabular}{|l|c|}
\hline
\textbf{Metrik} & \textbf{Nilai} \\
\hline
Jumlah Cluster & [isi setelah eksekusi] \\
Jumlah Noise Points & [isi setelah eksekusi] \\
Persentase Noise & [isi setelah eksekusi]\% \\
\hline
\end{tabular}
\end{table}

\subsection{Evaluasi Model}

\begin{table}[H]
\centering
\caption{Metrik Evaluasi Clustering}
\begin{tabular}{|l|c|l|}
\hline
\textbf{Metrik} & \textbf{Nilai} & \textbf{Interpretasi} \\
\hline
Silhouette Score & [isi setelah eksekusi] & [Baik/Cukup/Kurang baik] \\
Davies-Bouldin Index & [isi setelah eksekusi] & [Baik/Cukup/Kurang baik] \\
Calinski-Harabasz Score & [isi setelah eksekusi] & [Baik/Cukup/Kurang baik] \\
\hline
\end{tabular}
\end{table}

\subsection{Pembahasan}

\subsubsection{Keunggulan DBSCAN untuk Data Gempa Bumi}
\begin{enumerate}
    \item \textbf{Deteksi Cluster Arbitrary Shape}: DBSCAN mampu mendeteksi zona seismik dengan bentuk tidak beraturan yang mengikuti pola patahan atau zona subduksi.
    
    \item \textbf{Identifikasi Outlier}: Algoritma dapat membedakan gempa yang terisolasi (noise) dari zona seismik aktif (cluster), memberikan insight tentang gempa yang tidak biasa.
    
    \item \textbf{Tidak Memerlukan Jumlah Cluster}: Berbeda dengan K-Means, DBSCAN dapat menemukan jumlah cluster secara otomatis berdasarkan densitas data.
    
    \item \textbf{Robust terhadap Noise}: Algoritma tidak terpengaruh oleh outlier dalam pembentukan cluster.
\end{enumerate}

\subsubsection{Implikasi Praktis}
Hasil clustering ini dapat digunakan untuk:
\begin{enumerate}
    \item \textbf{Mitigasi Bencana}: Identifikasi zona seismik aktif untuk perencanaan tata kota dan bangunan tahan gempa.
    
    \item \textbf{Early Warning System}: Pemahaman pola cluster membantu pengembangan sistem peringatan dini gempa.
    
    \item \textbf{Penelitian Geologi}: Analisis pola gempa memberikan insight tentang pergerakan lempeng tektonik.
    
    \item \textbf{Asuransi dan Manajemen Risiko}: Pemetaan zona risiko gempa untuk perhitungan premi asuransi.
\end{enumerate}

\section{Kesimpulan}

Berdasarkan analisis clustering gempa bumi global menggunakan DBSCAN, dapat disimpulkan bahwa:

\begin{enumerate}
    \item Algoritma DBSCAN berhasil mengidentifikasi zona seismik aktif di berbagai belahan dunia dengan membentuk cluster yang merepresentasikan pola geografis gempa bumi.
    
    \item Parameter optimal yang digunakan adalah eps = 0.15 dan min\_samples = 5, yang ditentukan berdasarkan K-distance graph.
    
    \item DBSCAN mampu membedakan gempa yang terkonsentrasi di zona seismik aktif (cluster) dari gempa yang terisolasi (noise).
    
    \item Cluster-cluster yang terbentuk sesuai dengan zona tektonik aktif yang dikenal seperti Ring of Fire, zona subduksi Indonesia, dan patahan-patahan utama dunia.
    
    \item Visualisasi geospasial menggunakan GeoDataFrame, Folium, dan Contextily memberikan representasi yang efektif untuk memahami pola spasial gempa bumi.
    
    \item Hasil clustering ini dapat diaplikasikan untuk mitigasi bencana, pengembangan early warning system, penelitian geologi, dan manajemen risiko gempa bumi.
\end{enumerate}

\section{Daftar Pustaka}

\begin{enumerate}
    \item Ester, M., Kriegel, H. P., Sander, J., \& Xu, X. (1996). \textit{A density-based algorithm for discovering clusters in large spatial databases with noise}. In Kdd (Vol. 96, No. 34, pp. 226-231).
    
    \item Scikit-learn Documentation. (2024). \textit{DBSCAN Clustering}. https://scikit-learn.org/stable/modules/clustering.html\#dbscan
    
    \item GeoPandas Documentation. (2024). \textit{Introduction to GeoPandas}. https://geopandas.org/
    
    \item Folium Documentation. (2024). \textit{Python Data, Leaflet.js Maps}. https://python-visualization.github.io/folium/
    
    \item Schubert, E., Sander, J., Ester, M., Kriegel, H. P., \& Xu, X. (2017). \textit{DBSCAN revisited, revisited: why and how you should (still) use DBSCAN}. ACM Transactions on Database Systems (TODS), 42(3), 1-21.
\end{enumerate}

\end{document}
